% =====================================================================
% Section 2: Related Work
% =====================================================================
%
% Status: Draft §2 Related Work v1 (2026-05-16)
%
% Structure: three-paragraph form (PropRAG/NeocorRAG/Relink-style)
%   §2.1 Graph-based Retrieval for Multi-hop QA
%   §2.2 KG Reasoning and Query-time Chain Mining
%   §2.3 Personalized PageRank and Coverage-aware Reranking
%
% Each paragraph ends with an "In contrast, \methodname{} ..." sentence
% that positions our differences without overclaiming the prior work's
% weaknesses.
%
% Dependencies (BibTeX):
%   - angeli2015leveraging (already in main paper)
%   - gutierrez2025hipporag (already in main paper; HippoRAG 2)
%   - sun2024tog (already in main paper)
%   - huang2026relink (already in main paper)
%   - peng2026neocorrag (already in main paper)
%   - jimenez2024hipporag (HippoRAG; needs adding)
%   - you2024proprag (PropRAG; needs adding)
%   - edge2024graphrag (Microsoft GraphRAG; needs adding)
%   - haveliwala2002topic (Topic-Sensitive PageRank; needs adding)
%   - carbonell1998mmr (MMR; needs adding)
%   - lin2011submodular (Submodular IR; needs adding)
%   - trivedi2023ircot (IRCoT; needs adding)
%
% =====================================================================

\section{Related Work}
\label{sec:related}

\paragraph{Graph-based retrieval for multi-hop QA.}
Standard retrieval-augmented generation methods rank passages by dense
or lexical similarity to the question
\citep{karpukhin2020dpr,robertson2009bm25}, which is effective for
single-hop questions but tends to miss bridge passages whose surface
similarity to the question is low. Recent work augments this pipeline
with explicit graph structure to improve multi-hop coverage.
HippoRAG \citep{jimenez2024hipporag} and HippoRAG 2
\citep{gutierrez2025hipporag} extract OpenIE triples from the corpus
and run personalized PageRank over a corpus-wide graph of passage and
entity nodes seeded from the question. PropRAG \citep{you2024proprag}
replaces triples with propositions and discovers multi-hop evidence by
beam search over a proposition graph, motivated by the observation
that triples can lose context. GraphRAG \citep{edge2024graphrag} builds
community-summary graphs over the corpus and supports retrieval through
hierarchical summary aggregation, while hypergraph variants such as
HyperGraphRAG aggregate evidence across multi-granular community
structures. These systems share the strategy of complementing dense
retrieval with offline graph structure, but the propagation or search
is typically performed at the level of the full graph or its summaries.
\methodname{} adopts the same offline OpenIE-based indexing strategy
\citep{angeli2015leveraging,gutierrez2025hipporag}, and differs in that
it propagates evidence flow over a \emph{query-local} subgraph and adds
a retrieval-enhancement step that targets two failure modes of vanilla
local PPR: bridge passages that lie far from any question anchor, and
top rankings that collapse around a single entity neighborhood.

\paragraph{KG reasoning and query-time chain mining.}
A second line of work uses extracted relations as a reasoning substrate
that is traversed or updated at query time. ToG \citep{sun2024tog}
performs LLM-guided traversal over a pre-built knowledge graph,
selecting paths step by step, while related KG-RAG approaches
\citep{baek2023kaping} retrieve subgraphs and pass them to a language
model as context. Relink \citep{huang2026relink} addresses KG
incompleteness by dynamically reconstructing query-specific evidence
graphs, repairing missing facts as needed. NeocorRAG
\citep{peng2026neocorrag} keeps documents as the reader unit but
invokes a language model on retrieved text to mine entity-relation
chains as prompt augmentation. Interleaved chain-of-thought retrieval
methods such as IRCoT \citep{trivedi2023ircot} and Self-Ask
\citep{press2023selfask} similarly drive retrieval through
language-model-generated reasoning traces, alternating between
generation and lookup. In contrast, \methodname{} keeps the OpenIE
structure on the indexing side: the offline graph is neither traversed
by an LLM at query time nor regenerated per question, and chain-aware
retrieval is computed by graph propagation rather than by query-time
chain generation over retrieved text. The reader consumes
natural-language passages, and any query-time language model use is
restricted to extracting question-side anchors and requirements.

\paragraph{Personalized PageRank and coverage-aware reranking.}
\methodname{}'s retrieval and enhancement components draw on
long-standing techniques in information retrieval. Personalized
PageRank \citep{haveliwala2002topic,jeh2003scaling} provides a
principled way to bias graph propagation toward a query, and has been
used in graph-augmented retrieval to score nodes given a seed
distribution. On the reranking side, MMR \citep{carbonell1998mmr}
introduced diversity-aware reordering through marginal relevance
against an already-selected set, and submodular set functions
\citep{lin2011submodular} formalize coverage and diminishing returns for
selecting summary or evidence sets. Neural rerankers such as monoT5
\citep{nogueira2020monot5} produce strong per-passage scores via dense
cross-encoding, but operate on passages independently and do not model
coverage across the selected set. \methodname{} adapts both
graph-propagation and coverage-aware perspectives to the multi-hop
retrieval setting: propagation is restricted to a query-induced
subgraph and runs over OpenIE fact units rather than over passage or
word nodes directly, and the coverage signal is computed under a
noisy-OR model over question-side entity-relation requirements rather
than over query tokens or word-level features. These choices align the
retrieval and reranking objectives with the question's relational
structure rather than with surface similarity alone.
